\documentclass[11pt]{article}
\usepackage[margin=1in]{geometry}
\usepackage{hyperref}
\usepackage{listings}
\usepackage{xcolor}

\lstset{
  basicstyle=\ttfamily\small,
  breaklines=true,
  columns=fullflexible,
  frame=single,
  backgroundcolor=\color{gray!10}
}

\title{grep-dap: Study Phase Log}
\author{Claude (under direction of P. Cornillon)}
\date{2026-05-16}

\begin{document}
\maketitle

\section*{Purpose}

This log records the commands, queries, and decisions made while carrying
out the \emph{study phase} of the \texttt{grep-dap} project. The goal of
the study phase is to develop working knowledge of OPeNDAP (the Data
Access Protocol) and of the \texttt{pydap} client, with emphasis on
metadata, data retrieval, and subsetting. The substantive findings are
collected separately in
\texttt{docs/opendap\_readme.tex}; this file documents the process.

\section{Prompt \#1 (2026-05-16)}

\subsection*{Instruction}

\emph{``Read this document. Access the documentation for OPeNDAP and pydap
and summarize your findings in the \texttt{opendap\_readme.tex} file.
Focus on metadata, data retrieval, and subsetting.''}

\subsection*{Plan}

\begin{enumerate}
  \item Read \texttt{CLAUDE.md} to absorb the project framing (the
    OPeNDAP metadata problem; the GenAI/contrastive-learning angle;
    rules about Python/conda/matplotlib).
  \item Read \texttt{prompts/study\_phase.md} to confirm task scope.
  \item Create \texttt{docs/} and seed this log.
  \item Pull the OPeNDAP/DAP documentation that the project file
    points to (NASA Earthdata OPeNDAP portal) and the canonical
    OPeNDAP and DAP specification material.
  \item Pull the \texttt{pydap} user-facing documentation (intro,
    client usage, server, handlers).
  \item Distill into \texttt{opendap\_readme.tex} with the three
    requested focus areas as the structural backbone.
  \item Verify the produced LaTeX compiles cleanly.
\end{enumerate}

\subsection*{Thoughts before research}

A few observations that shaped the approach:

\begin{itemize}
  \item OPeNDAP, as a protocol, is older than its current web
    presence. The Hyrax server documentation and the OPeNDAP user
    guides are the most authoritative descriptions of DAP2 and DAP4
    semantics. For \texttt{grep-dap} the protocol details that
    matter are: what \emph{shapes} of responses a server can return
    (DDS / DAS / DDX / DMR / DAP), and what URL grammar a client
    uses to request subsets. Without those, ``infer the contents of
    an OPeNDAP server'' has no foothold.
  \item For \texttt{pydap}, what matters is (a) how a dataset is
    represented as a lazy in-Python tree, (b) how to inspect
    metadata without pulling data, and (c) how to express subsets
    so that only the bytes you actually want come down the wire.
  \item DAP4 introduces a different metadata document (DMR) and a
    different on-wire format. Code that ``infers contents'' must
    decide whether to lean on DAP2 (universal, simpler) or DAP4
    (richer, but not every server speaks it). Worth flagging in the
    readme.
\end{itemize}

\subsection*{Commands and tool calls}

\begin{lstlisting}
# Read project context
Read CLAUDE.md
Read prompts/study_phase.md

# Inspect existing docs directory
Glob docs/**/*   -> (empty)

# Create docs/ and seed this log file
Write docs/study_phase_log.tex

# Pull OPeNDAP / DAP documentation
WebFetch https://www.earthdata.nasa.gov/engage/open-data-services-software/earthdata-developer-portal/opendap
WebFetch https://www.opendap.org/about
WebSearch "OPeNDAP DAP2 DDS DAS constraint expression"
WebSearch "OPeNDAP DAP4 DMR DSR specification"

# Pull pydap documentation
WebFetch https://pydap.github.io/pydap/en/intro.html
WebFetch https://pydap.github.io/pydap/en/client.html
WebFetch (additional pydap pages as needed)
\end{lstlisting}

\subsection*{Notes during research}

The detailed findings live in \texttt{opendap\_readme.tex}.  Items
worth recording here as \emph{process} rather than \emph{content}:

\begin{itemize}
  \item The NASA Earthdata page is a portal/index; it points at
    Hyrax, at the DAP2/DAP4 specs, and at client tools, but it is
    not itself a specification. Treat as a pointer page.
  \item Some pydap doc pages are dynamic; if a fetch returns
    a near-empty shell, switch to the Chrome browsing tool so the
    page renders before extraction. (Not invoked unless needed.)
  \item Where the OPeNDAP and pydap docs disagree (e.g.\ DAP2 vs
    DAP4 default in newer pydap releases), follow pydap's own
    statement, since \texttt{grep-dap} will speak through pydap.
\end{itemize}

\subsection*{Outcome}

\begin{itemize}
  \item \texttt{docs/study\_phase\_log.tex} (this file) initialized.
  \item \texttt{docs/opendap\_readme.tex} written, structured as
    \emph{Overview}, \emph{Metadata}, \emph{Data retrieval},
    \emph{Subsetting}, \emph{pydap usage}, \emph{Implications for
    grep-dap}.
  \item Both files compile under \texttt{pdflatex}.
\end{itemize}

\end{document}
